% Fixed columns keep the table stable as rows appear.
% Arguments: vertical position, work, paradigm, task, soundness error.
\newcommand{\surveyrow}[5]{%
  \node[tablecell] at (-6.45,#1) {#2};
  \node[tablecell] at (-2.75,#1) {#3};
  \node[tablecell] at ( 1.05,#1) {#4};
  \node[tablecell] at ( 3.40,#1) {#5};
}

% outline.md, lines 60--75. References link directly to the papers.
\begin{frame}{Literature Survey}
  \centering
  \begin{tikzpicture}[
      tablecell/.style={anchor=west,inner sep=0pt,font=\small,text=CVQCInk},
      tablehead/.style={tablecell,font=\small\bfseries},
      separator/.style={draw=CVQCMuted!25,line width=0.4pt}
    ]
    \path[use as bounding box] (-6.7,-3.0) rectangle (6.7,2.9);

    \node[tablehead] at (-6.45,2.45) {Work};
    \node[tablehead] at (-2.75,2.45) {Paradigm};
    \node[tablehead] at ( 1.05,2.45) {Problem type};
    \node[tablehead] at ( 3.40,2.45) {Soundness error};
    \draw[draw=CVQCInk,line width=0.8pt] (-6.65,2.10) -- (6.65,2.10);

    \surveyrow{1.65}%
      {[\href{https://arxiv.org/abs/1904.06320}{GV19},
        \href{https://drops.dagstuhl.de/entities/document/10.4230/LIPIcs.ICALP.2023.67}{GMP23},
        \href{https://arxiv.org/abs/2609.38060}{HN26}]}%
      {Prepare-and-send}{Decision}{Negligible}

    \only<2->{
      \draw[separator] (-6.65,1.23) -- (6.65,1.23);
      \surveyrow{0.80}%
        {[\href{https://arxiv.org/abs/1804.01082}{Mah18},
          \href{https://arxiv.org/abs/1912.00990}{CCY20}]}%
        {Receive-and-measure}{Decision}{Negligible}
    }

    \only<3->{
      \draw[separator] (-6.65,0.38) -- (6.65,0.38);
      \surveyrow{-0.05}%
        {[\href{https://arxiv.org/abs/2012.04848}{CLLW22}]}%
        {Receive-and-measure}{Sampling}{Inverse-polynomial}
    }

    \only<4->{
      \fill[CVQCPale,rounded corners=3pt] (-6.65,-1.28) rectangle (6.65,-0.52);
      \begin{scope}[tablecell/.append style={font=\small\bfseries,text=CVQCTeal}]
        \surveyrow{-0.90}{This work}{Prepare-and-send}{Sampling}{Negligible}
      \end{scope}
    }

    \only<5>{
      \surveyrow{-1.75}%
        {[\href{https://arxiv.org/abs/2610.31775}{KKR26}]}%
        {Receive-and-measure}{Sampling}{Negligible}
    }

  \end{tikzpicture}

  \note<1>{
    \footnotesize
    Representative results, grouped by the two paradigms rather than an
    exhaustive chronology. GV19 is Alexandru Gheorghiu and Thomas Vidick,
    Computationally-secure and composable remote state preparation (FOCS 2019).
    GMP23 is Alexandru Gheorghiu, Tony Metger, and Alexander Poremba,
    Quantum Cryptography with Classical Communication (ICALP 2023).
    HN26 is Finn Holler and Anand Natarajan, Classical Verification of Quantum
    Computation with Quasilinear Resources, from Compiled Nonlocal Games (2026).
    They use a computational self-test and remote preparation to compile
    Broadbent's protocol, achieving quasilinear total resources under
    subexponential LWE hardness. Their stated BQP theorem has a constant gap.
    All three use remote preparation to remove quantum messages from the verifier.
    The negligible-error entry refers to decision verification after
    amplification. The preparation primitives have their own approximation
    guarantees; this row does not claim negligible-error remote preparation.
  }
  \note<2>{
    \footnotesize
    Mah18 is Urmila Mahadev, Classical Verification of Quantum Computations
    (FOCS 2018). CCY20 is Nai-Hui Chia, Kai-Min Chung, and Takashi Yamakawa,
    Classical Verification of Quantum Computations with Efficient Verifier
    (TCC 2020). The latter proves negligible soundness for parallel repetition
    of Mahadev's protocol. The table records the amplified decision guarantee,
    not the soundness of a single run of the original protocol.
  }
  \note<3>{
    \footnotesize
    CLLW22 is Kai-Min Chung, Yi Lee, Han-Hsuan Lin, and Xiaodi Wu,
    Constant-Round Blind Classical Verification of Quantum Sampling
    (EUROCRYPT 2022). It extends the Mahadev approach to sampling, but with
    inverse-polynomial soundness error. A majority vote that amplifies a
    decision bit does not generically turn several imperfect samples into a
    sample from the desired distribution. Soundness here concerns the whole
    sample-or-reject output distribution.
  }
  \note<4>{
    \footnotesize
    Our work follows the prepare-and-send route and achieves negligible
    soundness error for sampling. Emphasize the difference from the preceding
    row: the problem type stays sampling, while the soundness error becomes
    negligible. The next slide will separate our weak-quantum-verifier result
    from the fully classical compilation and state their assumptions.
  }
  \note<5>{
    \footnotesize
    KKR26 is Yael Tauman Kalai, Dakshita Khurana, and Justin Raizes,
    \href{https://arxiv.org/abs/2610.31775}{The Entire Quantum Computer on a
    Classical Leash (2026)}. This is concurrent and independent work, identified
    in the talk outline as STOC submission 598. Their framework builds on
    committed history states and non-destructive testing, and obtains
    negligible-error sampling verification under post-quantum LWE. The
    receive-and-measure label describes that underlying route; their full
    construction provides a more general classical--quantum two-party
    computation functionality.
  }
\end{frame}
